%%% Внесите свои данные - Input your data
%%
%%
\newcommand{\Author}{И.Р.\,Фёдорова} % И.О. Фамилия автора
\newcommand{\AuthorFull}{Фёдорова Ирина Романовна} % Фамилия Имя Отчество автора
\newcommand{\Supervisor}{В.А.\,Пархоменко} % И. О. Фамилия научного руководителя
\newcommand{\SupervisorJob}{старший преподаватель ВШ ПИ} %

\newcommand{\thesisYear}{2026} % заменить на год защиты
%%
%%
\newcommand{\group}{з5130903/40002} % заменить вместо N номер группы
\newcommand{\thesisSpecialtyCode}{09.03.03}% код направления подготовки
\newcommand{\thesisSpecialtyTitle}{Прикладная информатика} % наименование направления/специальности
%%
%%
\newcommand{\institute}{Институт компьютерных наук и~кибербезопасности}

\newcommand{\keywordsRu}{BDD, Gherkin, Cucumber.js, алгоритм Манакера, палиндромные подстроки, Node.js, нагрузочное тестирование} % ВВЕДИТЕ ключевые слова по-русски
\newcommand{\abstractRu}{В работе рассматривается применение BDD-подхода к консольному приложению на JavaScript для поиска палиндромных подстрок алгоритмом Манакера. Описаны основы Gherkin, выбор Cucumber.js, разработка пользовательских историй и реализация исполняемых BDD-сценариев. Для наукоемкой задачи выполнено сравнение двух алгоритмов: алгоритма Манакера и расширения от центра. Подготовлены генераторы датасетов, проведено профилирование времени выполнения и сформированы выводы о соответствии экспериментальных результатов теоретической сложности алгоритмов.}
%

\newcommand{\PracticeType}{Отчет о выполнении практической работы}
\newcommand{\School}{Высшая школа программной инженерии}
\newcommand{\practiceTitle}{Разработка кода через тестирование поведения}
\newcommand{\disciplineName}{Методы тестирования программного обеспечения} % название дисциплины


%%% Не меняем дальнейшую часть - Do not modify the rest part
%%
%%
%%
%%
\ifnumequal{\value{docType}}{1}{% Если ВКР, то...
	\newcommand{\DocType}{Выпускная квалификационная работа}
	\newcommand{\pdfDocType}{\DocType~(\thesisDegree)} %задаём метаданные pdf файла
	\newcommand{\pdfTitle}{\thesisTitle}
}{% Иначе 
	\newcommand{\DocType}{\PracticeType}
	\newcommand{\pdfDocType}{\DocType} %задаём метаданные pdf файла
	\newcommand{\pdfTitle}{\practiceTitle}
}%

\newcommand{\thesisSpecialtyCodeAndTitle}{\thesisSpecialtyCode~\thesisSpecialtyTitle}% Код и наименование направления/специальности
%%
%%
\hypersetup{%часть болка hypesetup в style
		pdftitle={\pdfTitle},    % Заголовок pdf-файла
		pdfauthor={\AuthorFull},    % Автор
		pdfsubject={\pdfDocType. Шифр и наименование направления подготовки: \thesisSpecialtyCodeAndTitle. \abstractRu},      % Тема
		pdfcreator={LaTeX, SPbPU-student-thesis-template},     % Приложение-создатель
%		pdfproducer={},  % Производитель, Производитель PDF % будет выставлена автоматически
		pdfkeywords={\keywordsRu}
}
%%
%%
%% вспомогательные команды
\newcommand{\firef}[1]{рис.\ref{#1}} %figure reference
\newcommand{\taref}[1]{табл.\ref{#1}}	%table reference
%%
%%


%%% Переопределение именований %%% Не меняем - Do not modify
\newcommand{\Ministry}{Министерство науки и высшего образования Российской~Федерации} %с 2020
\newcommand{\SPbPU}{Санкт-Петербургский политехнический университет Петра~Великого}
\newcommand{\SPbPUOfficialPrefix}{Федеральное государственное автономное образовательное учреждение высшего образования}
\newcommand{\SPbPUOfficialShort}{ФГАОУ~ВО~<<СПбПУ>>}
\renewcommand{\alsoname}{см. также}
\renewcommand{\seename}{см.}
\renewcommand{\headtoname}{вх.}
\renewcommand{\ccname}{исх.}
\renewcommand{\enclname}{вкл.}
\renewcommand{\pagename}{Pages}
\renewcommand{\partname}{Часть}
\renewcommand{\abstractname}{\textbf{Аннотация}}
\newcommand{\abstractnameENG}{\textbf{Annotation}}
\newcommand{\keywords}{\textbf{Ключевые слова}}
\newcommand{\keywordsENG}{\textbf{Keywords}}
\newcommand{\acknowledgements}{\textbf{Благодарности}}
\newcommand{\acknowledgementsENG}{\textbf{Acknowledgements}}
\renewcommand{\contentsname}{Content} %
\renewcommand{\figurename}{Рис.} % Стиль СПбПУ
\renewcommand{\tablename}{Таблица} % (ГОСТ Р 7.0.11-2011, 5.3.10)
\renewcommand{\listfigurename}{Список рисунков}
\renewcommand{\listtablename}{Список таблиц}
\renewcommand{\refname}{\fullbibtitle}
\renewcommand{\bibname}{\fullbibtitle}

\newcommand{\chapterEnTitle}{Сhapter title} % <- input the English title here (only once!) 
\newcommand{\chapterRuTitle}{Название главы}          % <- введите 
\newcommand{\sectionEnTitle}{Section title} %<- input subparagraph title in english
\newcommand{\sectionRuTitle}{Название подраздела} % <- введите название подраздела по-русски
\newcommand{\subsectionEnTitle}{Subsection title} % - input subsection title in english
\newcommand{\subsectionRuTitle}{Название параграфа} % <- введите название параграфа по-русски
\newcommand{\subsubsectionEnTitle}{Subsubsection title} % <- input subparagraph title in english
\newcommand{\subsubsectionRuTitle}{Название подпараграфа} % <- введите название подпараграфа по-русски
